\ifdefined\gkver\else\def\gkver{student}\fi
\documentclass[\gkver]{gaokaozhenti}
\begin{document}

\chapter{2008年江苏}

\begin{enumerate}
\item
%% number: 1
%% paperName: 2008年江苏高考第 1 题 3 分
%% typeId: 1
%% type: 单选题
%% chapter: 曲线运动
%% point: 万有引力定律
%% method: 无
%% score: 3
%% degree: 650
%% duplicateId: 0
%% body:
火星的质量和半径约为地球的 $\frac{1}{10}$ 和 $\frac{1}{2}$，地球表面的重力加速度为 $g$，则火星表面的重力加速度为\xzanswer{B}

\fourchoices[answer=B]
{$0.2g$}
{$0.4g$}
{$2.5g$}
{$5g$}

%% answer: B
%% memo:
\memoanswer{考查万有引力定律。星球表面重力等于万有引力，$G\frac{Mm}{R^{2}}=mg$，故火星表面的重力加速度 $g_{\text{火}}=\frac{M_{\text{火}}R_{\text{地}}^{2}}{M_{\text{地}}R_{\text{火}}^{2}}g=\frac{1/10}{(1/2)^{2}}g=0.4g$，故 B 正确。}

\item
%% number: 2
%% paperName: 2008年江苏高考第 2 题 3 分
%% typeId: 1
%% type: 单选题
%% chapter: 电路
%% point: 欧姆定律
%% method: 无
%% score: 3
%% degree: 850
%% duplicateId: 0
%% body:
2007 年度诺贝尔物理学奖授予了法国和德国的两位科学家，以表彰他们发现“巨磁电阻效应”，基于巨磁电阻效应开发的用于读取硬盘数据的技术被认为是纳米技术的第一次真正应用。在下列有关其他电阻应用的说法中，正确的是\xzanswer{D}

\fourchoices[answer=D]
{热敏电阻可用于温度测控装置中}
{光敏电阻是一种光电传感器}
{电阻丝可用于电热设备中}
{电阻在电路中主要起到通过直流、阻碍交流的作用}

%% answer: D
%% memo:
\memoanswer{考查基本物理常识。热敏电阻的原理是通过已知某电阻的电阻值与温度的函数关系，测得该热敏电阻的值即可获取温度，从而应用于温度测控装置中，A 说法正确；光敏电阻是将光信号与电信号进行转换的传感器，B 说法正确；电阻丝通过电流会产生热效应，可应用于电热设备中，C 说法正确；电阻对直流和交流均起到阻碍的作用，D 说法错误。}

\item
%% number: 3
%% paperName: 2008年江苏高考第 3 题 3 分
%% typeId: 1
%% type: 单选题
%% chapter: 运动和力
%% point: 牛顿第二定律
%% method: 无
%% score: 3
%% degree: 650
%% duplicateId: 0
%% body:
一质量为 $M$ 的探空气球在匀速下降，若气球所受浮力 $F$ 始终保持不变，气球在运动过程中所受阻力仅与速率有关，重力加速度为 $g$，现欲使该气球以同样的速率匀速上升，则需从气球吊篮中减少的质量为\xzanswer{A}

\onepicture[w=2.8cm]{figs/03.png}

\fourchoices[answer=A]
{$2\left(M-\frac{F}{g}\right)$}
{$M-2\frac{F}{g}$}
{$2M-\frac{F}{g}$}
{$0$}

%% answer: A
%% memo:
\memoanswer{考查牛顿运动定律。设减少的质量为 $\Delta m$，匀速下降时 $Mg=F+kv$，匀速上升时 $Mg-\Delta mg+kv=F$，解得 $\Delta m=2\left(M-\frac{F}{g}\right)$，A 正确。本题要注意受力分析各个力的方向。}

\item
%% number: 4
%% paperName: 2008年江苏高考第 4 题 3 分
%% typeId: 1
%% type: 单选题
%% chapter: 电路
%% point: 逻辑电路
%% method: 无
%% score: 3
%% degree: 850
%% duplicateId: 0
%% body:
在如图所示的逻辑电路中，当 A 端输入电信号“1”，B 端输入电信号“0”时，则在 C 和 D 端输出的电信号分别为\xzanswer{C}

\onepicture[w=5.0cm]{figs/04.png}

\fourchoices[answer=C]
{1 和 0}
{0 和 1}
{1 和 1}
{0 和 0}

%% answer: C
%% memo:
\memoanswer{正确认识门电路的符号，“\&”为或门，“1”为非门，其真值为：B 端 0 输入，则 1 输出，或门为“0，1”输入，则 1 输出。C 正确。}

\item
%% number: 5
%% paperName: 2008年江苏高考第 5 题 3 分
%% typeId: 1
%% type: 单选题
%% chapter: 直线运动
%% point: 运动图象
%% method: 无
%% score: 3
%% degree: 450
%% duplicateId: 0
%% body:
如图所示，粗糙的斜面与光滑的水平面相连接，滑块沿水平面以速度 $v_{0}$ 运动，设滑块运动的 A 点的时刻为 $t=0$，距 A 点的水平距离为 $x$，水平速度为 $v_{x}$。由于 $v_{0}$ 不同，从 A 点到 B 点的几种可能的运动图像如下列选项所示，其中表示摩擦力做功最大的是\xzanswer{D}

\onepicture[w=4.5cm]{figs/05a.png}

\onepicture[w=13.0cm]{figs/05b.png}

%% answer: D
%% memo:
\memoanswer{考查平抛运动的分解与牛顿运动定律。从 A 选项的水平位移与时间的正比关系可知，滑块做平抛运动，摩擦力必定为零；B 选项先平抛后在水平地面运动，水平速度突然增大，摩擦力依然为零；对 C 选项，水平速度不变，为平抛运动，摩擦力为零；对 D 选项水平速度与时间成正比，说明滑块在斜面上做匀加速直线运动，有摩擦力，故摩擦力做功最大的是 D 图像所显示的情景，D 对。本题考查非常灵活，但考查内容非常基础，抓住水平位移与水平速度与时间的关系，然后与平抛运动的思想结合起来，是为破解点。}

\item
%% number: 6
%% paperName: 2008年江苏高考第 6 题 4 分
%% typeId: 2
%% type: 多选题
%% chapter: 电场
%% point: 等势面
%% method: 无
%% score: 4
%% degree: 650
%% duplicateId: 0
%% body:
如图所示，实线为电场线，虚线为等势面。且 $AB=BC$，电场中 A、B、C 三点的场强分别为 $E_{A}$、$E_{B}$、$E_{C}$，电势分别为 $\varphi_{A}$、$\varphi_{B}$、$\varphi_{C}$，AB、BC 间的电势差分别为 $U_{AB}$、$U_{BC}$，则下列关系式中正确的是\xzanswer{ABC}

\onepicture[w=5.5cm]{figs/06.png}

\fourchoices[answer=ABC]
{$\varphi_{A}>\varphi_{B}>\varphi_{C}$}
{$E_{C}>E_{B}>E_{A}$}
{$U_{AB}<U_{BC}$}
{$U_{AB}=U_{BC}$}

%% answer: ABC
%% memo:
\memoanswer{考查静电场中的电场线、等势面的分布知识和规律。A、B、C 三点处在一根电场线上，沿着电场线的方向电势降落，故 $\varphi_{A}>\varphi_{B}>\varphi_{C}$，A 正确；由电场线的密集程度可看出电场强度大小关系为 $E_{C}>E_{B}>E_{A}$，B 对；电场线密集的地方电势降落较快，故 $U_{BC}>U_{AB}$，C 对 D 错。此类问题要在平时注重对电场线与场强、等势面与场强和电场线的关系的掌握，熟练理解常见电场线和等势面的分布规律。}

\item
%% number: 7
%% paperName: 2008年江苏高考第 7 题 4 分
%% typeId: 2
%% type: 多选题
%% chapter: 功和能
%% point: 机械能守恒定律
%% method: 无
%% score: 4
%% degree: 650
%% duplicateId: 0
%% body:
如图所示，两光滑斜面的倾角分别为 $30^{\circ}$ 和 $45^{\circ}$，质量分别为 $2m$ 和 $m$ 的两个滑块用不可伸长的轻绳跨过滑轮连接（不计滑轮的质量和摩擦），分别置于两个斜面上并由静止释放；若交换位置，再由静止释放。则在上述两种情况中正确的有\xzanswer{BD}

\onepicture[w=5.0cm]{figs/07.png}

\fourchoices[answer=BD]
{质量为 $2m$ 的滑块受到重力、绳的拉力、沿斜面的下滑力和斜面的支持力的作用}
{质量为 $m$ 的滑块均沿斜面向上运动}
{绳对质量为 $m$ 的滑块的拉力均大于该滑块对绳的拉力}
{系统在运动中机械能均守恒}

%% answer: BD
%% memo:
\memoanswer{考查受力分析、连接体整体法处理复杂问题的能力。每个滑块受到三个力：重力、绳子拉力、斜面的支持力，受力分析中应该是按性质分类的力，沿着斜面下滑力是分解出来的按照效果命名的力，A 错；对 B 选项，物体是上滑还是下滑要看两个物体的重力沿着斜面向下的分量的大小关系，由于 $2m$ 质量的滑块的重力沿着斜面的下滑分力较大，故质量为 $m$ 的滑块必定沿着斜面向上运动，B 对；任何一个滑块受到的绳子拉力与绳子对滑块的拉力等大反向，C 错；对系统除了重力之外，支持力对系统每个滑块不做功，绳子拉力对每个滑块的拉力等大反向，且对滑块的位移必定大小相等，故绳子拉力作为系统的内力对系统做功总和必定为零，故只有重力做功的系统，机械能守恒，D 对。}

\item
%% number: 8
%% paperName: 2008年江苏高考第 8 题 4 分
%% typeId: 2
%% type: 多选题
%% chapter: 电磁感应
%% point: 自感与互感，涡流
%% method: 无
%% score: 4
%% degree: 650
%% duplicateId: 0
%% body:
如图所示的电路中，三个相同的灯泡 a、b、c 和电感 $L_{1}$、$L_{2}$ 与直流电源连接，电感的电阻忽略不计，电键 K 从闭合状态突然断开时，下列判断正确的有\xzanswer{AD}

\onepicture[w=5.5cm]{figs/08.png}

\fourchoices[answer=AD]
{a 先变亮，然后逐渐变暗}
{b 先变亮，然后逐渐变暗}
{c 先变亮，然后逐渐变暗}
{b、c 都逐渐变暗}

%% answer: AD
%% memo:
\memoanswer{考查自感现象。电键 $K$ 闭合时，电感 $L_{1}$ 和 $L_{2}$ 的电流均等于三个灯泡的电流，断开电键 $K$ 的瞬间，电感上的电流 $i$ 突然减小，三个灯泡均处于回路中，故 $b$、$c$ 灯泡由电流 $i$ 逐渐减小，B、C 均错，D 对；原来每个电感线圈产生感应电动势均加载于灯泡 $a$ 上，故灯泡 $a$ 先变亮，然后逐渐变暗，A 对。本题涉及到自感现象中的“亮一下”现象，平时要注意透彻理解。}

\item
%% number: 9
%% paperName: 2008年江苏高考第 9 题 4 分
%% typeId: 2
%% type: 多选题
%% chapter: 曲线运动
%% point: 竖直平面圆周运动
%% method: 无
%% score: 4
%% degree: 450
%% duplicateId: 0
%% body:
如图所示，一根不可伸长的轻绳两端各系一个小球，跨在两根固定在同一高度的光滑水平细杆上，质量为 $3m$ 的 a 球置于地面上，质量为 $m$ 的 b 球从水平位置静止释放，当 a 球刚好对地面压力为零时，b 球摆过的角度为 $\theta$，下列结论正确的是\xzanswer{AC}

\onepicture[w=4.5cm]{figs/09.png}

\fourchoices[answer=AC]
{$\theta=90^{\circ}$}
{$\theta=45^{\circ}$}
{b 球摆动到最低点的过程中，重力对小球做功的功率先增大后减小}
{b 球摆动到最低点的过程中，重力对小球做功的功率一直增大}

%% answer: AC
%% memo:
\memoanswer{考查向心加速度公式、动能定理、功率等概念和规律。设 $b$ 球的摆动半径为 $R$，当摆过角度 $\theta$ 时的速度为 $v$，对 $b$ 球由动能定理 $mgR\sin\theta=\frac{1}{2}mv^{2}$，此时绳子拉力为 $T=3mg$，在绳子方向由向心力公式 $T-mg\sin\theta=m\frac{v^{2}}{R}$，解得 $\theta=90^{\circ}$，A 对 B 错；$b$ 球摆动到最低点的过程中一直机械能守恒，竖直方向的分速度先从零开始逐渐增大，然后逐渐减小到零，故重力的瞬时功率 $P_{b}=mgv_{\text{竖}}$ 先增大后减小，C 对 D 错。}

\item
%% number: 10
%% paperName: 2008年江苏高考第 10 题 8 分
%% typeId: 4
%% type: 实验题
%% chapter: 电路
%% point: 电阻定律
%% method: 无
%% score: 8
%% degree: 650
%% duplicateId: 0
%% body:
某同学想要了解导线在质量相同时，电阻与截面积的关系。选取了材料相同、质量相同的 5 卷导线，进行了如下实验：

\begin{enumerate}
	\item
	用螺旋测微器测量某一导线的直径如下图所示，读得直径 $d=$ \tkanswer[1.8cm]{} $\Umm$；
	
	\item
	该同学经实验测量及相关计算得到如下数据：
	
	\begin{table}[htbp]
	\centering
	\begin{tblr}{|c|c|c|c|c|c|}
	\hline
	电阻 $R$（$\UO$） & 121.0 & 50.0 & 23.9 & 10.0 & 3.1 \\
	\hline
	导线直径 $d$（$\Umm$） & 0.801 & 0.999 & 1.201 & 1.494 & 1.998 \\
	\hline
	导线截面积 $S$（$\Ummq$） & 0.504 & 0.784 & 1.133 & 1.753 & 3.135 \\
	\hline
	\end{tblr}
	\end{table}
	
	请你根据以上数据判断，该种导线的电阻 $R$ 与截面积 $S$ 是否满足反比关系？若满足反比关系请说明理由；若不满足，请写出 $R$ 与 $S$ 应满足的关系。
	
	\item
	若导线的电阻率 $\rho=5.1\times10^{-7}\UOm$，则表中阻值为 $3.1\UO$ 的导线长度 $L=$ \tkanswer[2cm]{} $\Um$（结果保留两位有效数字）。
\end{enumerate}

\onepicture[w=6.5cm, align=right]{figs/10.png}

%% answer:
\jdanswer{
\begin{enumerate}
	\item
	$1.200$
	
	\item
	不满足。由数据可知 $R$ 与 $S$ 不成反比，而 $R$ 与 $S^{2}$ 成反比，即 $RS^{2}$ 近似为常量。
	
	\item
	$19$
\end{enumerate}
}

%% memo:
\memoanswer{
（1）外径千分尺读数先读可动刻度左边界露出的主刻度 $L$，如本题为 $1\Umm$，再看主尺水平线对应的可动刻度 $n$，如本题为 20.0，记数为 $L+n\times0.01=1+20.0\times0.01=1.200\Umm$。读 $L$ 时要注意半毫米刻度线是否露出，如露出则记为 1.5 或 2.5 等。

（2）直接用两组 $R$、$S$ 值相乘（$50\times0.784=39.2$，$10.0\times1.753=17.53$），可得 $RS$ 明显不相等，可迅速判断结果“不满足”；同时可简单计算 $50.0\times0.999^{4}\approx50\times1$，$10\times1.494^{4}\approx10\times1.5^{4}=50$，两者接近相等，即 $R$ 与 $d$ 的四次方成反比，可迅速得出 $R$ 与 $S^{2}$ 成反比。计算时选择合适的数据可使计算简单方便，如本题中的（50.0，0.999，0.784）和（10.0，1.494，1.753）。

（3）根据 $R=\rho\frac{L}{S}$ 有 $L=\frac{RS}{\rho}=\frac{3.1\times3.135\times10^{-6}}{5.1\times10^{-7}}\Um\approx19\Um$。
}

\item
%% number: 11
%% paperName: 2008年江苏高考第 11 题 11 分
%% typeId: 4
%% type: 实验题
%% chapter: 曲线运动
%% point: 研究平抛运动实验
%% method: 无
%% score: 11
%% degree: 650
%% duplicateId: 0
%% body:
某同学利用如图所示的实验装置验证机械能守恒定律，弧形轨道末端水平，离地面的高度为 $H$，将钢球从轨道的不同高度 $h$ 处静止释放，钢球的落点距轨道末端的水平距离为 $s$。

\onepicture[w=7.0cm]{figs/11a.png}

\begin{enumerate}
	\item
	若轨道完全光滑，$s^{2}$ 与 $h$ 的理论关系应满足 $s^{2}=$ \tkanswer[2.5cm]{}（用 $H$、$h$ 表示）；
	
	\item
	该同学经实验测量得到一组数据，如下表所示：
	
	\begin{table}[htbp]
	\centering
	\begin{tblr}{|c|c|c|c|c|c|}
	\hline
	$h$（$10^{-1}\Um$） & 2.00 & 3.00 & 4.00 & 5.00 & 6.00 \\
	\hline
	$s^{2}$（$10^{-1}\Umq$） & 2.62 & 3.89 & 5.20 & 6.53 & 7.78 \\
	\hline
	\end{tblr}
	\end{table}
	
	请在坐标纸上作出 $s^{2}-h$ 关系图；
	
	\onepicture[w=8.0cm]{figs/11b.png}
	
	\item
	对比实验结果与理论计算得到的 $s^{2}-h$ 关系图线（图中已画出），自同一高度静止释放的钢球，水平抛出的速率 \tkanswer[1.5cm]{}（填“小于”或“大于”）理论值；
	
	\item
	从 $s^{2}-h$ 关系图中分析得出钢球水平抛出的速率差十分显著，你认为造成上述偏差的可能原因是 \tkanswer[4cm]{}。
\end{enumerate}

%% answer:
\jdanswer{
\begin{enumerate}
	\item
	$4Hh$
	
	\item
	图略（描点连线，不要画成折线）。
	
	\item
	小于
	
	\item
	轨道与钢球间存在摩擦，机械能减小；钢球的转动也需要能量维持。
\end{enumerate}
}

%% memo:
\memoanswer{
（1）根据机械能守恒，可得离开轨道时速度为 $\sqrt{2gh}$，由平抛运动知识可求得时间为 $\sqrt{\frac{2H}{g}}$，可得 $s=vt=\sqrt{4Hh}$。

（2）依次描点、连线，注意不要画成折线。

（3）从图中看，同一 $h$ 下的 $s^{2}$ 值，理论值明显大于实际值，而在同一高度 $H$ 下的平抛运动水平射程由水平速率决定，可见实际水平速率小于理论速率。

（4）由于客观上轨道与小球间存在摩擦，机械能减小，因此会导致实际值比理论值小；小球的转动也需要能量维持，而机械能守恒中没有考虑重力势能转化成转动能的这一部分，也会导致实际速率明显小于“理论”速率。
}

\item
%% number: 12
%% paperName: 2008年江苏高考第 12 题 12 分
%% typeId: 3
%% type: 填空题
%% chapter: 原子和原子核
%% point: 半衰期
%% method: 无
%% score: 12
%% degree: 650
%% duplicateId: 0
%% body:
选做题：请从 A、B 和 C 三小题中选定两小题作答，并在答题卡上把所选题目对应字母后的方框涂满涂黑。如都作答，按 A、B 两小题评分。

\textbf{A．}（选修模块 3-3）（12 分）

\begin{enumerate}
	\item
	空气压缩机在一次压缩过程中，活塞对气缸中气体做功为 $2.0\times10^{5}\UJ$，同时气体的内能增加了 $1.5\times10^{5}\UJ$，气体 \tkanswer[1.5cm]{}（填“吸收”或“放出”）的热量等于 \tkanswer[2cm]{} $\UJ$。
	
	\item
	若一定质量的理想气体分别按下图所示的三种不同过程变化，其中表示等压变化的是 \tkanswer[1.5cm]{}（填“A”、“B”或“C”），该过程中气体的内能 \tkanswer[1.5cm]{}（填“增加”、“减少”或“不变”）。
	
	\item
	设想将 $1\Ug$ 水均匀分布在地球表面上，估算 $1\Ucmq$ 表面上有多少个水分子？（已知 $1\Umol$ 水的质量为 $18\Ug$，地球的表面积约为 $5\times10^{14}\Umq$，结果保留一位有效数字）
\end{enumerate}

\onepicture[w=11.0cm]{figs/12A.png}

\textbf{B．}（选修模块 3-4）（12 分）

\begin{enumerate}
	\item
	一列沿着 $x$ 轴正方向传播的横波，在 $t=0$ 时刻的波形如图甲所示，图甲中某质点的振动图像如图乙所示，质点 N 的振幅是 \tkanswer[1.5cm]{} $\Um$，振动周期为 \tkanswer[1.2cm]{} $\Us$，图乙表示质点 \tkanswer[1.2cm]{}（从质点 K、L、M、N 中选填）的振动图像，该波的波速为 \tkanswer[1.5cm]{} $\Ums$。
	
	\onepicture[w=10.0cm]{figs/12Ba.png}
	
	\item
	惯性系 S 中有一边长为 $l$ 的正方形（如图 A 所示），从相对 S 系沿 $x$ 轴方向以接近光速匀速飞行的飞行器上测得该正方形的图像是 \tkanswer[1.5cm]{}。
	
	\onepicture[w=12.0cm]{figs/12Bb.png}
	
	\item
	描述简谐运动特征的公式是 $x=$ \tkanswer[2.5cm]{}，自由下落的篮球经地面反弹后上升又落下，若不考虑空气阻力及在地面反弹时的能量损失，此运动 \tkanswer[1.5cm]{}（填“是”或“不是”）简谐运动。
\end{enumerate}

\textbf{C．}（选修模块 3-5）（12 分）

\begin{enumerate}
	\item
	下列实验中，深入地揭示了光的粒子性一面的有 \tkanswer[2cm]{}。
	
	（A）X 射线被石墨散射后部分波长增大
	
	（B）锌板被紫外线照射时有电子逸出，但被可见光照射时没有电子逸出
	
	（C）轰击金箔的 $\alpha$ 粒子有少数运动方向发生较大偏转
	
	（D）光量子发射的光经三棱镜分光后出现线状光谱
	
	\onepicture[w=13.0cm]{figs/12Ca.png}
	
	\item
	场强为 $E$、方向竖直向上的匀强电场中的两小球 A、B，它们的质量分别为 $m_{1}$、$m_{2}$，电量分别为 $q_{1}$、$q_{2}$，A、B 两球由静止释放，重力加速度为 $g$，则小球 A、B 组成的系统动量守恒应满足的关系式为 \tkanswer[5cm]{}。
	
	\item
	约里奥·居里夫妇因发现人工放射性而获得了 1935 年的诺贝尔化学奖，他们发现的放射性元素 $\ce{^{30}_{15}P}$ 衰变成 $\ce{^{30}_{14}Si}$ 的同时放出另一种粒子，这种粒子是 \tkanswer[2cm]{}，$\ce{^{32}_{15}P}$ 是 $\ce{^{30}_{15}P}$ 的同位素，被广泛应用于生物示踪技术，$1\Umg$ $\ce{^{32}_{15}P}$ 随时间衰变的关系如图所示，请估算 $4\Umg$ 的 $\ce{^{32}_{15}P}$ 经多少天的衰变后还剩 $0.25\Umg$？
	
	\onepicture[w=8.0cm]{figs/12Cb.png}
\end{enumerate}

%% answer:
\jdanswer{
\begin{enumerate}
	\item
	A：（1）放出，$5\times10^{4}\UJ$；（2）C，增加；（3）$n=\frac{m}{M}N_{A}\frac{S}{S_{0}}\approx7\times10^{3}$ 个。
	
	\item
	B：（1）$0.8\Um$，$4\Us$，L，$0.5\Ums$；（2）C；（3）$x=A\sin\omega t$，不是。
	
	\item
	C：（1）A、B；（2）$E(q_{1}+q_{2})=(m_{1}+m_{2})g$；（3）正电子，56 天。
\end{enumerate}
}

%% memo:
\memoanswer{
\textbf{A．}（1）由热力学第一定律 $\Delta U=W+Q$，代入数据得 $1.5\times10^{5}=2.0\times10^{5}+Q$，解得 $Q=-5\times10^{4}\UJ$，即气体放出 $5\times10^{4}\UJ$ 的热量。（2）由 $\frac{PV}{T}=$ 恒量，压强不变时 $V$ 随温度 $T$ 的变化是一次函数关系，故选择 C 图；该过程温度升高，内能增加。（3）$1\Ug$ 水的分子数 $N=\frac{m}{M}N_{A}$，$1\Ucmq$ 表面上的分子数 $n=N\frac{S}{S_{0}}\approx7\times10^{3}$（$6\times10^{3}\sim7\times10^{3}$ 都算对）。

\textbf{B．}（1）从甲、乙图可看出波长 $\lambda=2.0\Um$，周期 $T=4\Us$，振幅 $A=0.8\Um$；乙图中显示 $t=0$ 时刻该质点处于平衡位置向上振动，甲图波形图中波向 $x$ 轴正方向传播，则 L 质点正在平衡位置向上振动，波速 $v=\frac{\lambda}{T}=0.5\Ums$。（2）由相对论知识易得运动方向上的边长变短，垂直运动方向的边长不变，C 图像正确。（3）简谐运动的特征公式为 $x=A\sin\omega t$，其中 $A$ 是振幅；自由下落的篮球由反弹起来的过程中，回复力始终为重力，恒定不变，与偏离平衡位置的位移不是成正比的，不符合简谐运动的规律。

\textbf{C．}（1）A 为康普顿散射，B 为光电效应，康普顿散射和光电效应都深入揭示了光的粒子性；C 为 $\alpha$ 粒子散射，不是光子，揭示了原子的核式结构模型；D 为光的折射，揭示了氢原子能级的不连续。（2）系统动量守恒的条件为所受合外力为零，即电场力与重力平衡，$E(q_{1}+q_{2})=(m_{1}+m_{2})g$。（3）由核反应过程中电荷数和质量数守恒可写出核反应方程 $\ce{^{30}_{15}P -> ^{30}_{14}Si + ^{0}_{1}e}$，可知这种粒子是正电子。由图像可知 $\ce{^{32}_{15}P}$ 的半衰期为 14 天，$4\Umg$ 的 $\ce{^{32}_{15}P}$ 衰变后还剩 $0.25\Umg$，经历了 4 个半衰期，所以为 56 天。
}

\item
%% number: 13
%% paperName: 2008年江苏高考第 13 题 15 分
%% typeId: 6
%% type: 计算题
%% chapter: 曲线运动
%% point: 平抛运动
%% method: 无
%% score: 15
%% degree: 450
%% duplicateId: 0
%% body:
抛体运动在各类体育运动项目中很常见，如乒乓球运动。现讨论乒乓球发球问题，设球台长 $2L$、网高 $h$，乒乓球反弹前后水平分速度不变，竖直分速度大小不变、方向相反，且不考虑乒乓球的旋转和空气阻力。（设重力加速度为 $g$）

\begin{enumerate}
	\item
	若球在球台边缘 O 点正上方高度为 $h_{1}$ 处以速度 $v_{1}$ 水平发出，落在球台的 $P_{1}$ 点（如图实线所示），求 $P_{1}$ 点距 O 点的距离 $x_{1}$；
	
	\item
	若球在 O 点正上方以速度 $v_{2}$ 水平发出，恰好在最高点时越过球网，落在球台的 $P_{2}$ 点（如图虚线所示），求 $v_{2}$ 的大小；
	
	\item
	若球在 O 点正上方水平发出后，球经反弹恰好越过球网且刚好落在对方球台边缘 $P_{3}$ 处，求发球点距 O 点的高度 $h_{3}$。
\end{enumerate}

\onepicture[w=6.5cm, align=right]{figs/13.png}

%% answer:
\jdanswer{
\begin{enumerate}
	\item
	$x_{1}=v_{1}\sqrt{\frac{2h_{1}}{g}}$
	
	\item
	$v_{2}=\frac{L}{2}\sqrt{\frac{g}{2h}}$
	
	\item
	$h_{3}=\frac{4}{3}h$
\end{enumerate}
}

%% memo:
\memoanswer{
（1）据平抛规律得 $h_{1}=\frac{1}{2}gt_{1}^{2}$，$x_{1}=v_{1}t_{1}$，解得 $x_{1}=v_{1}\sqrt{\frac{2h_{1}}{g}}$。

（2）同理得 $h_{2}=\frac{1}{2}gt_{2}^{2}$，$x_{2}=v_{2}t_{2}$，且 $h_{2}=h$，$2x_{2}=L$，解得 $v_{2}=\frac{L}{2}\sqrt{\frac{g}{2h}}$。

（3）如图，同理得 $h_{3}=\frac{1}{2}gt_{3}^{2}$，$x_{3}=v_{3}t_{3}$，且 $3x_{3}=2L$。设球从恰好越过球网到最高点的时间为 $t$，水平距离为 $s$，有 $h_{3}-h=\frac{1}{2}gt^{2}$，$s=v_{3}t$，由几何关系得 $x_{3}+s=L$，解得 $h_{3}=\frac{4}{3}h$。

\onepicture[w=7.5cm]{figs/13_an.jpg}
}

\item
%% number: 14
%% paperName: 2008年江苏高考第 14 题 16 分
%% typeId: 6
%% type: 计算题
%% chapter: 磁场
%% point: 带电粒子在磁场中的运动
%% method: 无
%% score: 16
%% degree: 650
%% duplicateId: 0
%% body:
在磁感应强度为 $B$ 的匀强磁场中，一质量为 $m$、带正电 $q$ 的小球在 O 点静止释放，小球的运动曲线如图所示，已知此曲线在最低点的曲率半径为该点到 $x$ 轴距离的 2 倍，重力加速度为 $g$。求：

\begin{enumerate}
	\item
	小球运动到任意位置 $P(x,y)$ 处的速率 $v$；
	
	\item
	小球在运动过程第一次下降的最大距离 $y_{m}$；
	
	\item
	当在上述磁场中加一竖直向上、场强为 $E$（$E>\frac{mg}{q}$）的匀强电场时，小球从 O 点静止释放后获得的最大速率 $v_{m}$。
\end{enumerate}

\onepicture[w=6.0cm, align=right]{figs/14.png}

%% answer:
\jdanswer{
\begin{enumerate}
	\item
	$v=\sqrt{2gy}$
	
	\item
	$y_{m}=\frac{2m^{2}g}{q^{2}B^{2}}$
	
	\item
	$v_{m}=\frac{2}{qB}(qE-mg)$
\end{enumerate}
}

%% memo:
\memoanswer{
（1）洛伦兹力不做功，由动能定理得 $mgy=\frac{1}{2}mv^{2}$ ①，解得 $v=\sqrt{2gy}$ ②。

（2）设在最大距离 $y_{m}$ 处的速率为 $v_{m}$，根据圆周运动有 $qBv_{m}-mg=m\frac{v_{m}^{2}}{R}$ ③，且由②得 $v_{m}=\sqrt{2gy_{m}}$ ④，由③④及 $R=2y_{m}$，得 $y_{m}=\frac{2m^{2}g}{q^{2}B^{2}}$ ⑤。

（3）小球运动如图所示，由动能定理得 $(qE-mg)|y_{m}|=\frac{1}{2}mv_{m}^{2}$ ⑥，由圆周运动得 $qBv_{m}+mg-qE=m\frac{v_{m}^{2}}{R}$ ⑦，由⑥⑦且 $R=2|y_{m}|$，解得 $v_{m}=\frac{2}{qB}(qE-mg)$。

\onepicture[w=6.5cm]{figs/14_an.jpg}
}

\item
%% number: 15
%% paperName: 2008年江苏高考第 15 题 16 分
%% typeId: 6
%% type: 计算题
%% chapter: 电磁感应
%% point: 电磁感应能量分析
%% method: 无
%% score: 16
%% degree: 450
%% duplicateId: 0
%% body:
如图所示，间距为 $l$ 的两条足够长的金属平行导轨与水平面的夹角为 $\theta$，导轨光滑且电阻不计，磁感应强度为 $B$ 的条形匀强磁场方向与导轨平面垂直，磁场区域的宽度为 $d_{1}$，间距为 $d_{2}$，两根质量均为 $m$、有效电阻均为 $R$ 的导体棒 a 和 b 放在导轨上，并与导轨垂直。（设重力加速度为 $g$）

\begin{enumerate}
	\item
	若 a 进入第二个磁场区域时，b 以与 a 同样的速率进入第一个磁场区域，求 b 穿过第一个磁场区域过程中增加的动能 $\Delta E_{k}$；
	
	\item
	若 a 进入第二个磁场区域时，b 恰好离开第一个磁场区域，此后 a 离开第二个磁场区域时，b 又恰好进入第二个磁场区域，且 a、b 在任意一个磁场区域或无磁场区域的运动时间相等，求 a 穿过第二个磁场区域过程中，两导体棒产生的总焦耳热 $Q$；
	
	\item
	对于第（2）问所述的运动情况，求 a 穿出第 $k$ 个磁场区域时的速率 $v_{k}$。
\end{enumerate}

\onepicture[w=8.0cm, align=right]{figs/15.png}

%% answer:
\jdanswer{
\begin{enumerate}
	\item
	$\Delta E_{k}=mgd_{1}\sin\theta$
	
	\item
	$Q=mg(d_{1}+d_{2})\sin\theta$
	
	\item
	$v_{k}=v_{1}=\frac{4mgRd_{2}\sin\theta}{B^{2}l^{2}d_{1}}-\frac{B^{2}l^{2}d_{1}}{8mR}$
\end{enumerate}
}

%% memo:
\memoanswer{
（1）a 和 b 不受安培力作用，由机械能守恒得 $\Delta E_{k}=mgd_{1}\sin\theta$ ①。

（2）设导体棒刚进入无磁场区域时的速度为 $v_{1}$，刚离开无磁场区域时的速度为 $v_{2}$，由能量守恒得：在磁场区域有 $\frac{1}{2}mv_{1}^{2}+Q=\frac{1}{2}mv_{2}^{2}+mgd_{1}\sin\theta$ ②；在无磁场区域有 $\frac{1}{2}mv_{2}^{2}=\frac{1}{2}mv_{1}^{2}+mgd_{2}\sin\theta$ ③。解得 $Q=mg(d_{1}+d_{2})\sin\theta$ ④。

（3）在无磁场区域时，根据匀变速直线运动规律有 $v_{2}-v_{1}=gt\sin\theta$ ⑤，且平均速度 $\frac{v_{1}+v_{2}}{2}=\frac{d_{2}}{t}$ ⑥；在有磁场区域，棒 a 受到的合力 $F=mg\sin\theta-BIl$ ⑦，感应电动势 $E=Blv$ ⑧，感应电流 $I=\frac{E}{2R}$ ⑨，解得 $F=mg\sin\theta-\frac{B^{2}l^{2}}{2R}v$ ⑩。根据牛顿第二定律，在 $t\sim t+\Delta t$ 时间内 $\Sigma\Delta v=\Sigma\frac{F}{m}\Delta t$（11），则有 $\Sigma\Delta v=\Sigma\left(g\sin\theta-\frac{B^{2}l^{2}v}{2mR}\right)\Delta t$（12），解得 $v_{2}-v_{1}=gt\sin\theta-\frac{B^{2}l^{2}}{2mR}d_{1}$（13）。联立⑤⑥（13）式解得 $v_{1}=\frac{4mgRd_{2}\sin\theta}{B^{2}l^{2}d_{1}}-\frac{B^{2}l^{2}d_{1}}{8mR}$，由题意知 $v=v_{1}$，即 $v_{k}=\frac{4mgRd_{2}\sin\theta}{B^{2}l^{2}d_{1}}-\frac{B^{2}l^{2}d_{1}}{8mR}$。
}

\end{enumerate}

\end{document}
